\section{Related Work}
\label{sec:related_work}

\subsection{Image-Conditioned 3D Generation}

Image-conditioned 3D generation has advanced through two complementary paradigms.
%
Feed-forward reconstruction methods directly infer 3D geometry and appearance from input images using learned geometric priors~\cite{ICLR2024_dcad3425_lrm,xu2024instantmesh,wang2024crm,boss2025sf3d}.
%
In contrast, generative approaches learn structured or native 3D latent representations to synthesize high-fidelity assets with complex geometry and appearance~\cite{Xiang_2025_CVPR_Trellis,Xiang_2026_CVPR_Trellis2,zhao2025hunyuan3d,li2025triposg,li2025step1x3d}.
%
Recent studies further improve fidelity to input observations through pixel-aligned conditioning and reconstruction-guided generation~\cite{li2026pixal3d,chang2026reconviagen,huang2026unirecgen}.
%
Among these, Pixal3D~\cite{li2026pixal3d} incorporates projected image features into 3D generation to better preserve observed appearance.
%
However, these methods primarily assume that the conditioning images depict the geometry to be generated.
%
In contrast, our framework combines a geometry-aligned target view with original source views depicting different geometric states, explicitly retargeting source appearance evidence across geometric variations.

























\subsection{Geometry-Guided 3D Editing}

Geometry-guided 3D editing aims to modify existing 3D assets while preserving their structural plausibility and visual cues.
%
Prior methods achieve geometric modifications through deformation constraints, learned shape priors, and image- or text-guided optimization~\cite{yoo2024plausible,choi2026image,kim2025meshup,dinh2025geometryinstyle}.
%
More recent approaches incorporate spatial controls into generative frameworks or introduce dedicated 3D editing models~\cite{fedele2026spacecontrol,erkoc2025preditor3d}.
%
However, geometric controls alone may leave the visual conditioning inconsistent with the edited shape.
%
Complementary to geometry editing, reference-guided texturing methods synthesize consistent appearance on prescribed geometry through diffusion-based view generation and texture transfer~\cite{chen2023text2tex,cao2023texfusion,yeh2024texturedreamer,xie2024styletex}.
%
Specifically, ShapeUP~\cite{gat2026shapeup} performs geometry editing through latent-space translation and subsequently guides texture synthesis using the edited geometry and renderings of the source asset.
%
While ShapeUP demonstrates the importance of preserving source appearance during geometric modifications, our approach uses a VLM-edited target view as a structural anchor and explicitly retargets source-view features across geometric states within the 3D generation process.





















\subsection{Correspondence-Guided Feature Fusion}

Dense correspondence estimation has benefited from semantic representations learned by diffusion and self-supervised vision models~\cite{tang2023emergent,zhang2023tale,luo2023diffusionhyperfeatures,simeoni2025dinov3}.
%
Beyond semantic descriptors, learned matching frameworks and geometric foundation models provide complementary structural and 3D-aware priors for establishing correspondences~\cite{Edstedt_2024_CVPR_roma,mast3r_eccv24,wang2025vggt,wang2026vggtomega}.
%
In particular, RoMa-$\Omega$~\cite{nordstrom2026romaomega} investigates the correspondence capabilities of VGGT-$\Omega$~\cite{wang2026vggtomega} features and incorporates selected geometric representations into a learned dense matching framework.
%
Its findings suggest that geometrically informative features do not necessarily provide reliable correspondences through direct feature similarity, highlighting the importance of dedicated feature matching.

Correspondence and geometric alignment have also been explored for consistent multi-view generation.
%
Existing methods employ synchronized denoising, cross-view attention, and geometry-aware feature aggregation to maintain consistency across generated views~\cite{shi2024mvdream,liu2024syncdreamer,hollein2024viewdiff,tang2023mvdiffusion,kant2024spad,huang2025mv,zhang2026align}.
%
While these approaches primarily focus on correspondence estimation or cross-view consistency, our work leverages semantic correspondences for a different purpose: retrieving source-specific appearance information across geometric state changes and adapting the matched features to guide target-state 3D generation.
